%%%%%%%%%%%%%%%%%%%%%%%%%%%%%%%%%%%%%%%%%%%%%%%%%%%%%%%%%%%%%%%%%%%%%%%%%%%%%%%%%%
\begin{frame}[fragile]\frametitle{}
\begin{center}
{\Large Implementation using Langchain}
\end{center}
\end{frame}

%%%%%%%%%%%%%%%%%%%%%%%%%%%%%%%%%%%%%%%%%%%%%%%%%%%%%%%%%%%
\begin{frame}[fragile]\frametitle{What is Graph RAG?}
  \begin{itemize}
    \item Enhances traditional RAG by incorporating knowledge graphs.
    \item Moves beyond vector similarity to structured entity relationships.
    \item Enables LLMs to reason over graph-based data structures.
  \end{itemize}
\end{frame}

%%%%%%%%%%%%%%%%%%%%%%%%%%%%%%%%%%%%%%%%%%%%%%%%%%%%%%%%%%%
\begin{frame}[fragile]\frametitle{Installation}
  \begin{itemize}
    \item Install required packages:
  \end{itemize}
  \begin{lstlisting}[language=bash]
pip install --upgrade --quiet \
  langchain langchain-community \
  langchain-core langchain-experimental \
  networkx json-repair
  \end{lstlisting}
\end{frame}

%%%%%%%%%%%%%%%%%%%%%%%%%%%%%%%%%%%%%%%%%%%%%%%%%%%%%%%%%%%
\begin{frame}[fragile]\frametitle{Importing Modules}
  \begin{itemize}
    \item Import essential modules:
  \end{itemize}
  \begin{lstlisting}[language=python]
from langchain_experimental.graph_transformers import LLMGraphTransformer
from langchain_experimental.graphs import GraphQAChain
from langchain_community.graphs.networkx_graph import NetworkxEntityGraph
from langchain_core.documents import Document
from langchain_google_genai import GoogleGenerativeAI
import google.generativeai as genai
  \end{lstlisting}
\end{frame}

%%%%%%%%%%%%%%%%%%%%%%%%%%%%%%%%%%%%%%%%%%%%%%%%%%%%%%%%%%%
\begin{frame}[fragile]\frametitle{Setting Up the LLM}
  \begin{itemize}
    \item Initialize the LLM (e.g., Google Generative AI):
  \end{itemize}
  \begin{lstlisting}[language=python]
  GOOGLE_API_KEY = ''
  genai.configure(api_key=GOOGLE_API_KEY)
llm = GoogleGenerativeAI(model="gemini-pro",api_key=GOOGLE_API_KEY)
  \end{lstlisting}
\end{frame}

%%%%%%%%%%%%%%%%%%%%%%%%%%%%%%%%%%%%%%%%%%%%%%%%%%%%%%%%%%%
\begin{frame}[fragile]\frametitle{Preparing the Text}
  \begin{itemize}
    \item Load and prepare the input text:
  \end{itemize}
  \begin{lstlisting}[language=python]
text = "Marie Curie was a Polish and naturalized-French physicist and chemist..."
doc = Document(page_content=text)  # LLMGraphTransformer requires Document objects
  \end{lstlisting}
\end{frame}

%%%%%%%%%%%%%%%%%%%%%%%%%%%%%%%%%%%%%%%%%%%%%%%%%%%%%%%%%%%
\begin{frame}[fragile]\frametitle{Transforming Text to Graph Documents}
  \begin{itemize}
    \item Convert text into graph documents:
  \end{itemize}
  \begin{lstlisting}[language=python]
transformer = LLMGraphTransformer(llm=llm)
graph_documents = transformer.convert_to_graph_documents([doc])
  \end{lstlisting}
\end{frame}

%%%%%%%%%%%%%%%%%%%%%%%%%%%%%%%%%%%%%%%%%%%%%%%%%%%%%%%%%%%
\begin{frame}[fragile]\frametitle{Specifying Entities and Relationships}
  \begin{itemize}
    \item Define allowed node types and relationship types; pass them to the constructor:
  \end{itemize}
  \begin{lstlisting}[language=python]
allowed_nodes = ["Person", "Country", "Organization"]
allowed_relationships = ["nationality", "located_in", "worked_at", "spouse", "mother"]
# Pass allowed_nodes/relationships in constructor, not convert method
transformer_filtered = LLMGraphTransformer(
    llm=llm,
    allowed_nodes=allowed_nodes,
    allowed_relationships=allowed_relationships
)
graph_documents_filtered = transformer_filtered.convert_to_graph_documents([doc])
  \end{lstlisting}
\end{frame}

%%%%%%%%%%%%%%%%%%%%%%%%%%%%%%%%%%%%%%%%%%%%%%%%%%%%%%%%%%%
\begin{frame}[fragile]\frametitle{Creating the Graph Object}
  \begin{itemize}
    \item Initialize the graph and add nodes and edges:
  \end{itemize}
  \begin{lstlisting}[language=python]
graph = NetworkxEntityGraph()
for graph_doc in graph_documents_filtered:
    for node in graph_doc.nodes:
        graph.add_node(node)
    for edge in graph_doc.relationships:
        graph.add_edge(edge.source.id, edge.target.id, edge.type)
  \end{lstlisting}
\end{frame}

%%%%%%%%%%%%%%%%%%%%%%%%%%%%%%%%%%%%%%%%%%%%%%%%%%%%%%%%%%%
\begin{frame}[fragile]\frametitle{Performing Graph-Based QA}
  \begin{itemize}
    \item Initialize the Graph QA chain and ask questions:
  \end{itemize}
  \begin{lstlisting}[language=python]
qa_chain = GraphQAChain.from_llm(llm=llm, graph=graph)
response = qa_chain.run("Tell me about Marie Curie")
print(response)
  \end{lstlisting}
\end{frame}

%%%%%%%%%%%%%%%%%%%%%%%%%%%%%%%%%%%%%%%%%%%%%%%%%%%%%%%%%%%
\begin{frame}[fragile]\frametitle{Conclusion}
  \begin{itemize}
    \item Graph RAG integrates structured knowledge into LLM responses.
    \item Enhances the depth and accuracy of information retrieval.
    \item Suitable for complex, relationship-rich data domains.
  \end{itemize}
\end{frame}



